\subsection{Determinants}

\begin{definition}{Minor Submatrix}{minor-submatrix}
  For $\mathbf{A} \in \R^{n \times n}$, let $\mathbf{A}_{ij} \in \R^{(n-1) \times (n-1)}$ denote the submatrix obtained from $\mathbf{A}$ by removing the $i$-th row and $j$-th column.
\end{definition}

\begin{definition}{Determinant}{determinant}
  Let $\mathbf{A} \in \R^{n \times n}$. The \textbf{determinant} $\det(\mathbf{A}) \in \R$ is defined inductively:
  \begin{itemize}
    \item $n = 1$: $\det(\mathbf{A}) = a_{11}$.
    \item $n = 2$: $\det(\mathbf{A}) = a_{11} a_{22} - a_{12} a_{21}$.
    \item $n = 3$:
          \[
            \det(\mathbf{A})
            = a_{11} \det
            \begin{bmatrix}
              a_{22} & a_{23} \\ a_{32} & a_{33}
            \end{bmatrix}
            - a_{12} \det
            \begin{bmatrix}
              a_{21} & a_{23} \\ a_{31} & a_{33}
            \end{bmatrix}
            + a_{13} \det
            \begin{bmatrix}
              a_{21} & a_{22} \\ a_{31} & a_{32}
            \end{bmatrix}
            .
          \]
  \end{itemize}

  And for general $n$, expanding along the first row:
  \[
    \det(\mathbf{A})
    = a_{11} \det(\mathbf{A}_{11})
    - a_{12} \det(\mathbf{A}_{12})
    + a_{13} \det(\mathbf{A}_{13})
    - a_{14} \det(\mathbf{A}_{14})
    + \cdots
    + (-1)^{n+1} a_{1n} \det(\mathbf{A}_{1n})
  \]

\end{definition}

\subsubsection{Cofactor Expansion Along Any Row}

\begin{definition}{Cofactor}{cofactor}
  For $\mathbf{A} \in \R^{n \times n}$, the \textbf{$(i,j)$ cofactor} is
  \[
    C_{ij} = (-1)^{i+j} \det(\mathbf{A}_{ij}).
  \]
\end{definition}

\begin{theorem}{Cofactor Expansion}{cofactor-expansion}
  For any row $i$ of $\mathbf{A} \in \R^{n \times n}$,
  \[
    \det(\mathbf{A}) = a_{i1} C_{i1} + a_{i2} C_{i2} + \cdots + a_{in} C_{in}.
  \]
\end{theorem}

\begin{example}
  Let
  \[
    \mathbf{A} =
    \begin{bmatrix}
      2 & 3 & 1 & 4 & 5  \\
      0 & 3 & 1 & 0 & -2 \\
      0 & 0 & 2 & 0 & 5  \\
      0 & 0 & 0 & 4 & -1 \\
      0 & 0 & 0 & 0 & -1
    \end{bmatrix}
    .
  \]
  Expand along row $5$: only $a_{55} = -1$ is non-zero, so
  \[
    \det(\mathbf{A}) = a_{55} C_{55} = (-1) \det(\mathbf{A}_{55})
    = -\det
    \begin{bmatrix}
      2 & 3 & 1 & 4 \\
      0 & 3 & 1 & 0 \\
      0 & 0 & 2 & 0 \\
      0 & 0 & 0 & 4
    \end{bmatrix}
    .
  \]
  Repeatedly expanding along the last row of each successive submatrix,
  \[
    \det(\mathbf{A}) = (-1)(4)(2)(3)(2) = -48,
  \]
  the product of the diagonal entries.
\end{example}

\begin{theorem}{Determinant of a Triangular Matrix}{triangular-det}
  If $\mathbf{A} \in \R^{n \times n}$ is upper or lower triangular, then
  \[
    \det(\mathbf{A}) = a_{11} a_{22} \cdots a_{nn},
  \]
  the product of its diagonal entries.
\end{theorem}

\subsection{Properties of Determinants}

Let $\mathbf{A} \in \R^{n \times n}$ with rows $\mathbf{R}_1, \dots, \mathbf{R}_n$.

\begin{theorem}{Row Operations and the Determinant}{row-ops-det}
  \begin{enumerate}
    \item \textbf{Swap.} Swapping two rows flips the sign:
          \[
            \det
            \begin{bmatrix}
              \vdots \\ \mathbf{R}_i \\ \vdots \\ \mathbf{R}_j \\ \vdots
            \end{bmatrix}
            = -\det
            \begin{bmatrix}
              \vdots \\ \mathbf{R}_j \\ \vdots \\ \mathbf{R}_i \\ \vdots
            \end{bmatrix}
            .
          \]
    \item \textbf{Scale.} Multiplying one row by $\lambda$ multiplies the determinant by $\lambda$:
          \[
            \det
            \begin{bmatrix}
              \vdots \\ \lambda \mathbf{R}_i \\ \vdots
            \end{bmatrix}
            = \lambda \det
            \begin{bmatrix}
              \vdots \\ \mathbf{R}_i \\ \vdots
            \end{bmatrix}
            .
          \]
    \item \textbf{Replace.} Adding a multiple of one row to another leaves the determinant unchanged:
          \[
            \det
            \begin{bmatrix}
              \vdots \\ \mathbf{R}_i + \lambda \mathbf{R}_j \\ \vdots
            \end{bmatrix}
            = \det
            \begin{bmatrix}
              \vdots \\ \mathbf{R}_i \\ \vdots
            \end{bmatrix}
            .
          \]
  \end{enumerate}
\end{theorem}

\begin{example}
  Reduce to triangular form to evaluate
  \[
    \det
    \begin{bmatrix}
      0  & 1 & 2  & 1  \\
      2  & 1 & -1 & 3  \\
      1  & 3 & 0  & -1 \\
      -1 & 4 & 1  & 5
    \end{bmatrix}
    .
  \]
  Swap $\mathbf{R}_1 \leftrightarrow \mathbf{R}_3$:
  \[
    = -\det
    \begin{bmatrix}
      1  & 3 & 0  & -1 \\
      2  & 1 & -1 & 3  \\
      0  & 1 & 2  & 1  \\
      -1 & 4 & 1  & 5
    \end{bmatrix}
    .
  \]
  $\mathbf{R}_2 \to \mathbf{R}_2 - 2\mathbf{R}_1$, $\mathbf{R}_4 \to \mathbf{R}_4 + \mathbf{R}_1$:
  \[
    = -\det
    \begin{bmatrix}
      1 & 3  & 0  & -1 \\
      0 & -5 & -1 & 5  \\
      0 & 1  & 2  & 1  \\
      0 & 7  & 1  & 4
    \end{bmatrix}
    .
  \]
  Swap $\mathbf{R}_2 \leftrightarrow \mathbf{R}_3$:
  \[
    = \det
    \begin{bmatrix}
      1 & 3  & 0  & -1 \\
      0 & 1  & 2  & 1  \\
      0 & -5 & -1 & 5  \\
      0 & 7  & 1  & 4
    \end{bmatrix}
    .
  \]
  $\mathbf{R}_3 \to \mathbf{R}_3 + 5\mathbf{R}_2$, $\mathbf{R}_4 \to \mathbf{R}_4 - 7\mathbf{R}_2$, then $\mathbf{R}_4 \to \mathbf{R}_4 + \tfrac{13}{9}\mathbf{R}_3$:
  \[
    = \det
    \begin{bmatrix}
      1 & 3 & 0 & -1             \\
      0 & 1 & 2 & 1              \\
      0 & 0 & 9 & 10             \\
      0 & 0 & 0 & \tfrac{103}{9}
    \end{bmatrix}
    = (1)(1)(9)\!\left(\tfrac{103}{9}\right) = 103.
  \]
\end{example}

\subsubsection{Invertibility and Determinants}

\begin{theorem}{Invertibility Criterion}{det-invertible}
  $\mathbf{A} \in \R^{n \times n}$ is invertible $\iff \det(\mathbf{A}) \neq 0$.
\end{theorem}

\begin{definition}{Singular Matrix}{singular}
  $\mathbf{A} \in \R^{n \times n}$ is \textbf{singular} (non-invertible) if $\det(\mathbf{A}) = 0$.
\end{definition}

\begin{theorem}{Determinant and Linear Independence}{det-lin-indep}
  The columns of $\mathbf{A} \in \R^{n \times n}$ are linearly independent $\iff \det(\mathbf{A}) \neq 0$.
\end{theorem}

\begin{remark}
  Two immediate consequences:
  \begin{enumerate}
    \item If $\mathbf{A}$ has two identical rows (or columns), then $\det(\mathbf{A}) = 0$.
    \item If $\mathbf{A}$ has a zero row (or column), then $\det(\mathbf{A}) = 0$.
  \end{enumerate}
\end{remark}

\subsubsection{Multiplicativity and Linearity}

\begin{theorem}{Product Rule}{det-product}
  For $\mathbf{A}, \mathbf{B} \in \R^{n \times n}$,
  \[
    \det(\mathbf{A}\mathbf{B}) = \det(\mathbf{A}) \det(\mathbf{B}).
  \]
\end{theorem}

\begin{remark}
  $\det(\mathbf{A} + \mathbf{B})$ is \emph{not} generally equal to $\det(\mathbf{A}) + \det(\mathbf{B})$.
\end{remark}

\begin{example}
  Take $\mathbf{A} =
    \begin{bmatrix}
      1 & 2 \\ 3 & 4
    \end{bmatrix}
  $, $\mathbf{B} =
    \begin{bmatrix}
      -1 & -2 \\ -3 & -4
    \end{bmatrix}
  $. Then $\mathbf{A} + \mathbf{B} = \mathbf{0}$, so
  \[
    \det(\mathbf{A} + \mathbf{B}) = 0,
    \quad
    \det(\mathbf{A}) + \det(\mathbf{B}) = (-2) + (-2) = -4.
  \]
\end{example}

\begin{theorem}{Determinant is Linear in Each Column}{det-column-linear}
  Write $\mathbf{A} = [\mathbf{a}_1 \mid \cdots \mid \mathbf{a}_{i-1} \mid \mathbf{x} \mid \mathbf{a}_{i+1} \mid \cdots \mid \mathbf{a}_n]$ with the $i$-th column variable. Define $T : \R^n \to \R$ by
  \[
    T(\mathbf{x}) = \det[\mathbf{a}_1 \mid \cdots \mid \mathbf{x} \mid \cdots \mid \mathbf{a}_n].
  \]
  Then $T$ is a linear transformation.
\end{theorem}

\begin{example}
  Splitting the second column $
    \begin{bmatrix}
      2 \\ 3
    \end{bmatrix}
    =
    \begin{bmatrix}
      1 \\ 0
    \end{bmatrix}
    +
    \begin{bmatrix}
      1 \\ 3
    \end{bmatrix}
  $:
  \[
    \det
    \begin{bmatrix}
      2 & 2 \\ 3 & 3
    \end{bmatrix}
    = \det
    \begin{bmatrix}
      2 & 1 \\ 3 & 0
    \end{bmatrix}
    + \det
    \begin{bmatrix}
      2 & 1 \\ 3 & 3
    \end{bmatrix}
    = (-3) + 3 = 0.
  \]

\end{example}

\begin{example}
  In general, for $2 \times 2$ matrices,
  \begin{align*}
    \det
    \begin{bmatrix}
      a + a' & b + b' \\ c + c' & d + d'
    \end{bmatrix}
     & = \det
    \begin{bmatrix}
      a + a' & b \\ c + c' & d
    \end{bmatrix}
    + \det
    \begin{bmatrix}
      a + a' & b' \\ c + c' & d'
    \end{bmatrix}
    \\
     & = \det
    \begin{bmatrix}
      a & b \\ c & d
    \end{bmatrix}
    + \det
    \begin{bmatrix}
      a' & b \\ c' & d
    \end{bmatrix}
    + \det
    \begin{bmatrix}
      a & b' \\ c & d'
    \end{bmatrix}
    + \det
    \begin{bmatrix}
      a' & b' \\ c' & d'
    \end{bmatrix}
  \end{align*}
\end{example}

\subsubsection{Transpose and Cramer's Rule}

\begin{theorem}{Determinant of Transpose}{det-transpose}
  For $\mathbf{A} \in \R^{n \times n}$,
  \[
    \det(\mathbf{A}^T) = \det(\mathbf{A}).
  \]
\end{theorem}

A row operation on $\mathbf{A}$ is a column operation on $\mathbf{A}^T$, so column operations obey the same rules as row operations.

\begin{theorem}{Cramer's Rule}{cramers-rule}
  Let $\mathbf{A} \in \R^{n \times n}$ be invertible and $\mathbf{b} \in \R^n$. The unique solution to $\mathbf{A}\mathbf{x} = \mathbf{b}$ has components
  \[
    x_j = \frac{\det(\mathbf{A}_j(\mathbf{b}))}{\det(\mathbf{A})},
  \]
  where $\mathbf{A}_j(\mathbf{b})$ is the matrix obtained by replacing the $j$-th column of $\mathbf{A}$ with $\mathbf{b}$.
\end{theorem}

\begin{example}
  Solve $\mathbf{A}\mathbf{x} = \mathbf{b}$ for
  \[
    \mathbf{A} =
    \begin{bmatrix}
      2 & 3 \\ 4 & -1
    \end{bmatrix}
    ,\quad
    \mathbf{b} =
    \begin{bmatrix}
      1 \\ 5
    \end{bmatrix}
    .
  \]
  $\det(\mathbf{A}) = -2 - 12 = -14$, so
  \[
    x_1 = \frac{\det
      \begin{bmatrix}
        1 & 3 \\ 5 & -1
      \end{bmatrix}
    }{\det(\mathbf{A})}
    = \frac{-16}{-14} = \frac{8}{7},
    \quad
    x_2 = \frac{\det
      \begin{bmatrix}
        2 & 1 \\ 4 & 5
      \end{bmatrix}
    }{\det(\mathbf{A})}
    = \frac{6}{-14} = -\frac{3}{7}.
  \]
\end{example}

\subsubsection{Adjugate Formula for the Inverse}

\begin{definition}{Adjugate}{adjugate}
  For $\mathbf{A} \in \R^{n \times n}$, the \textbf{adjugate} of $\mathbf{A}$ is the transpose of the cofactor matrix:
  \[
    \operatorname{adj}(\mathbf{A}) = [C_{ij}]^T =
    \begin{bmatrix}
      C_{11} & C_{21} & \cdots & C_{n1} \\
      \vdots & \vdots & \ddots & \vdots \\
      C_{1n} & C_{2n} & \cdots & C_{nn}
    \end{bmatrix}
    ,
  \]
  where $C_{ij} = (-1)^{i+j} \det(\mathbf{A}_{ij})$ is the $(i,j)$ cofactor.
\end{definition}

\begin{theorem}{Adjugate Formula for the Inverse}{adj-inverse}
  If $\mathbf{A} \in \R^{n \times n}$ is invertible, then
  \[
    \mathbf{A}^{-1} = \frac{\operatorname{adj}(\mathbf{A})}{\det(\mathbf{A})}.
  \]
\end{theorem}

\begin{theorem}{Determinant of the Adjugate}{det-adjugate}
  For $\mathbf{A} \in \R^{n \times n}$ invertible,
  \[
    \det(\operatorname{adj}(\mathbf{A})) = (\det(\mathbf{A}))^{n-1}.
  \]
\end{theorem}

\begin{proof}
  Taking determinants of both sides of $\mathbf{A}^{-1} = \tfrac{1}{\det(\mathbf{A})} \operatorname{adj}(\mathbf{A})$,
  \[
    \det(\mathbf{A}^{-1})
    = \det\!\left(\tfrac{1}{\det(\mathbf{A})} \operatorname{adj}(\mathbf{A})\right).
  \]
  The left side equals $\tfrac{1}{\det(\mathbf{A})}$. On the right, scaling all $n$ columns by $\tfrac{1}{\det(\mathbf{A})}$ pulls out $\left(\tfrac{1}{\det(\mathbf{A})}\right)^n$ by linearity in each column:
  \[
    \frac{1}{\det(\mathbf{A})}
    = \left(\frac{1}{\det(\mathbf{A})}\right)^n \det(\operatorname{adj}(\mathbf{A})).
  \]
  Solving gives $\det(\operatorname{adj}(\mathbf{A})) = (\det(\mathbf{A}))^{n-1}$.
\end{proof}
